% Chapter 10: Stratified sampling and conditional Monte Carlo.

The methods of \cref{ch:08-antithetic-control} use a symmetry of the input law or a
quantity with a known mean. This chapter uses the structure of the variance itself. For any
random variable $\xi$ computed along with the payoff $Y$, the law of total variance splits
\begin{equation}\label{eq:st-total-variance}
  \Var Y=\underbrace{\E\big[\Var(Y\mid\xi)\big]}_{\text{within}}+\underbrace{\Var\big(\E[Y\mid\xi]\big)}_{\text{between}} .
\end{equation}
\emph{Stratified sampling} cuts the range of $\xi$ into pieces and fixes in advance how
many draws land in each piece; it removes the variance between the pieces, and in the limit
of many pieces the whole between part. \emph{Conditional Monte Carlo} computes $\E[Y\mid\xi]$ in closed form and
averages that instead of $Y$; it removes the within part. The two are complementary, and
the art is in the choice of $\xi$: the terminal value of a Brownian path, a direction in the
space of normals, every coordinate at once (Latin hypercube sampling), or the whole
volatility path of a stochastic volatility model, given which the Black--Scholes formula
applies.

Throughout, $Y$ is a real random variable with $\E[Y^2]<\infty$, mean $\mu$ and variance
$\sigma^2$, and $\hat\mu_n$ is the plain estimator from $n$ independent copies, with
variance $\sigma^2/n$ (\cref{def:mc-estimator}). In pricing, $Y$ is a discounted payoff and
the samples are drawn under $\Q$; as in \cref{ch:06-monte-carlo} we write $\Prob$ and $\E$
for the sampling measure. The test case is $S=K=100$, $T=1$, $r=1\%$, $y=2\%$,
$\sigma=20\%$, with $S_T=S\exp(\nu T+\sigma\sqrt T Z)$, $\nu=r-y-\frac{\sigma^2}2$, and the
discounted call has mean $7.36429$ and variance $158.006$.

\section{Stratified sampling}\label{sec:st-strata}

Let $\xi$ be a random element of a measurable space $E$ (the \emph{stratification
variable}) and let $A_1,\dots,A_H$ be a measurable partition of $E$ into \emph{strata}, with
$p_h=\Prob(\xi\in A_h)$. A stratum with $p_h=0$ is almost surely never visited and
contributes nothing to $\mu$; we discard such strata and assume $p_h>0$ for every $h$. Let
\[
  \mu_h=\E[Y\mid\xi\in A_h]=\frac{\E[Y\1_{A_h}(\xi)]}{p_h},\qquad
  \sigma_h^2=\E\big[(Y-\mu_h)^2\mid\xi\in A_h\big]
\]
be the mean and the variance of $Y$ in stratum $h$. Given integers $n_h\ge1$ with
$\sum_hn_h=n$, draw $Y_{h,1},\dots,Y_{h,n_h}$ independently from the conditional law of $Y$
given $\xi\in A_h$, independently across strata, and set
\begin{equation}\label{eq:st-estimator}
  \hat\mu^{\mathrm{st}}_n=\sum_{h=1}^Hp_h\,\bar Y_h,\qquad \bar Y_h=\frac1{n_h}\sum_{i=1}^{n_h}Y_{h,i}.
\end{equation}
Two things are needed: the probabilities $p_h$ exactly, and a way to sample in a stratum.
Both are easy when $\xi$ is a uniform. Suppose $Y=g(U,U')$, where $U$ is uniform on $(0,1)$
and independent of the rest $U'$ of the randomness, and take strata
$I_h=(a_{h-1},a_h]$ with $0=a_0<\dots<a_H=1$, so $p_h=a_h-a_{h-1}$. If $V$ is uniform and
independent of $U'$, then $g(a_{h-1}+p_hV,U')$ has the conditional law of $Y$ given
$U\in I_h$: for bounded measurable $F$, Fubini's theorem gives
$\E[F(Y)\1_{I_h}(U)]=\int_{I_h}\E[F(g(u,U'))]\,\dd u=p_h\E[F(g(a_{h-1}+p_hV,U'))]$. With
$Z=N^{-1}(U)$ (\cref{prop:rng-inversion}) this stratifies a normal; equiprobable strata
$a_h=h/H$ are the usual choice. Defining the strata in $u$-space makes every $p_h$ exact
even when the variable built from $U$ has atoms or gaps.

\begin{theorem}[Stratified sampling]\label{thm:st-variance}
In the setting above:
\begin{enumerate}[label=(\roman*)]
  \item $\mu=\sum_hp_h\mu_h$ and
    $\sigma^2=\sum_hp_h\sigma_h^2+\sum_hp_h(\mu_h-\mu)^2$ (within plus between).
  \item If $n_h\ge1$ for every $h$, $\hat\mu^{\mathrm{st}}_n$ is unbiased and
    $\Var\hat\mu^{\mathrm{st}}_n=\sum_hp_h^2\sigma_h^2/n_h$. If $n_h=0$ for the strata in a set
    $E_0$ and their terms are dropped, the bias is $-\sum_{h\in E_0}p_h\mu_h$.
  \item (\emph{Proportional allocation.}) If $n_h=np_h$ for every $h$ (integers), then
    \[
      \Var\hat\mu^{\mathrm{st}}_n=\frac1n\sum_hp_h\sigma_h^2=\frac{\sigma^2}n-\frac1n\sum_hp_h(\mu_h-\mu)^2\le\frac{\sigma^2}{n},
    \]
    with equality if and only if $\mu_h=\mu$ for every $h$.
  \item Splitting strata further never increases $\sum_hp_h\sigma_h^2$.
  \item (\emph{Many strata.}) Let $\xi=U$ be uniform and $I_h=((h-1)/H,h/H]$. As
    $H\to\infty$, $\sum_hp_h\sigma_h^2\to\E[\Var(Y\mid U)]=\sigma^2-\Var(\E[Y\mid U])$. In
    particular, if $Y$ is a function of $U$ alone and $H=H_n\to\infty$ with proportional
    allocation, then $n\Var\hat\mu^{\mathrm{st}}_n\to0$.
\end{enumerate}
\end{theorem}

\begin{proof}
(i) $\E Y=\sum_h\E[Y\1_{A_h}(\xi)]=\sum_hp_h\mu_h$. Expanding $Y-\mu=(Y-\mu_h)+(\mu_h-\mu)$ on
$\{\xi\in A_h\}$, the cross term has expectation
$2(\mu_h-\mu)\E[(Y-\mu_h)\1_{A_h}(\xi)]=0$, so
$\E[(Y-\mu)^2\1_{A_h}(\xi)]=p_h\sigma_h^2+p_h(\mu_h-\mu)^2$; sum over $h$.

(ii) The $\bar Y_h$ are independent, with means $\mu_h$ and variances $\sigma_h^2/n_h$; use
(i). Without draws in the strata of $E_0$ the estimator misses their terms $p_h\mu_h$ in (i).

(iii) Substitute $n_h=np_h$ in (ii) and use (i). The difference is
$\frac1n\sum_hp_h(\mu_h-\mu)^2$, which vanishes exactly when every $\mu_h$ equals $\mu$,
since every $p_h>0$.

(iv) If $A_h$ is split into pieces $A_{h,k}$, apply (i) to the conditional law given
$\xi\in A_h$: $\sigma_h^2$ is at least the within-piece part
$\sum_k(p_{h,k}/p_h)\sigma_{h,k}^2$. Multiply by $p_h$ and sum over $h$.

(v) Let $g(U)=\E[Y\mid U]$ and let $\mathcal G_H$ be generated by the events
$\{U\in I_h\}$. On $\{U\in I_h\}$, $\E[Y\mid\mathcal G_H]=\mu_h$, so
$\sum_hp_h\sigma_h^2=\E[(Y-\E[Y\mid\mathcal G_H])^2]$. Since $\mathcal G_H\subset\sigma(U)$,
$\E[Y\mid\mathcal G_H]=\E[g(U)\mid\mathcal G_H]=:(P_Hg)(U)$, where $P_Hg$ is the average of $g$
over each cell $I_h$. The variable $Y-g(U)$ is orthogonal to every square-integrable function of $U$, so
\[
  \sum_hp_h\sigma_h^2=\E\big[(Y-g(U))^2\big]+\|g-P_Hg\|^2_{L^2(0,1)} .
\]
The first term is $\E[\Var(Y\mid U)]$. For the second, $P_H$ is an orthogonal projection, so
$\|P_Hk\|\le\|k\|$; for a continuous $c$ on $[0,1]$, $|c-P_Hc|\le\omega_c(1/H)$, the modulus
of continuity of $c$ at $1/H$. Hence
$\|g-P_Hg\|\le\|g-c\|+\|c-P_Hc\|+\|P_H(c-g)\|\le2\|g-c\|+\omega_c(1/H)$. Continuous functions
are dense in $L^2(0,1)$, so the second term tends to $0$.
\end{proof}

Under plain sampling the number of draws that land in $A_h$ is random, binomial with
mean $np_h$; proportional allocation fixes it at its mean, and the between-strata variance,
which came from that randomness, disappears. The word ``never'' in (iii) belongs to
proportional allocation. With other allocations stratification can lose: take two
equiprobable strata with $\mu_1=\mu_2$, $\sigma_1=0$ and $\sigma_2>0$, so that
$\sigma^2=\sigma_2^2/2$; the allocation $n_1=n-1$, $n_2=1$ gives the variance
$\sigma_2^2/4>\sigma^2/n$ for every $n\ge3$. Part (v) says that with many strata on a uniform,
stratification removes all the variance that the uniform explains. By
\cref{prop:rng-inversion} the same holds for equiprobable strata of any scalar $\xi=F^{-1}(U)$
with a continuous distribution function $F$, since then $U=F(\xi)$ and $\sigma(\xi)=\sigma(U)$. No rate is
claimed: how fast $\sum_hp_h\sigma_h^2$ approaches its limit depends on the smoothness of
$u\mapsto\E[Y\mid U=u]$ (\exref{ex:st-one-point}).

\begin{example}[The test-case call]\label{exa:st-call}
Stratify $Z=N^{-1}(U)$ into $H$ equiprobable strata. The call pays off when
$Z>z^*=0.15$, so every stratum below $z^*$ has $\mu_h=\sigma_h=0$: with $H=2$ the lower
half is worthless, and the variance factor $\sum_hp_h\sigma_h^2/\sigma^2$ is $0.657$. With
$H=10$, $100$ and $1000$ it is $0.1007$, $0.0081$ and $0.00073$, standard-error ratios
$0.317$, $0.090$ and $0.027$ (\cref{fig:st-strata}). The stratum moments are computed by
Gauss--Legendre quadrature on each stratum (\texttt{stratum\_moments}), which avoids the
cancellation of the closed forms in narrow strata. A stratified draw costs one inversion. On
our machine \texttt{rng.norm\_inv} costs about seven times NumPy's normal generator, and the
whole stratified estimator about $4.5$ times the plain one with NumPy normals (machine-dependent;
\texttt{laboratory\_costs} and the notebook time it), so the efficiency
(\cref{def:mc-efficiency}) is about $1/(4.5\times0.1007)\approx2.2$ with $10$ strata, not $10$,
and about $27$ with $100$.
\end{example}

\begin{figure}[ht]
  \centering
  \includegraphics[width=\textwidth]{ch10-strata}
  \caption{(a) The discounted test-case call as a function of the uniform $u$,
    $g(u)=e^{-rT}(S_T(N^{-1}(u))-K)^+$, and its means on $8$ equiprobable strata; the curve
    leaves the frame near $u=1$. Proportional stratification samples only the deviations
    from the step function. (b) Variance factors $n\Var\hat\mu^{\mathrm{st}}_n/\sigma^2$
    against the number $H$ of equiprobable strata: lines are exact
    (\cref{thm:st-variance,prop:st-optimal}), dots are measured over $400$ independent
    replications of $n=2000$ draws, seed $1010$ (filled: proportional; open: Neyman, with at
    least one draw per stratum). The deep out-of-the-money call is the course's example.}
  \label{fig:st-strata}
\end{figure}

\paragraph{The course's experiment.} The course's notebook prices a deep
out-of-the-money call, $S=100$, $K=150$, $T=1$, $r=5\%$, $y=0$, $\sigma=20\%$ (price
$0.35963$), with $n=10^5$ draws in $10$, $50$ or $100$ equiprobable strata, and compares the
variances of $100$ repeated estimates. The exact variance factors are $0.856$, $0.318$ and
$0.158$. Only $3.0\%$ of the mass lies above the strike ($z^*=1.877$), so with $10$ strata
the whole payoff sits in the top stratum, and stratification removes only the randomness of
the number of draws that reach it; Neyman allocation repairs this (\exref{ex:st-neyman-otm}). Repeating the course's protocol in our implementation
(\texttt{course\_otm\_experiment}, numpy \texttt{default\_rng}, seed $5$), the variances of
the $100$ estimates are $9.61$, $8.05$, $2.64$ and $1.46$ (in units of $10^{-5}$) for plain
sampling and $10$, $50$, $100$ strata, against the exact $8.08$, $6.92$, $2.57$ and $1.28$; a
variance estimated from $100$ estimates has a relative standard error of about $14\%$.

\subsection{Optimal allocation}

Proportional allocation ignores the $\sigma_h$. Strata where the payoff varies more deserve
more draws.

\begin{proposition}[Neyman allocation]\label{prop:st-optimal}
Let $\bar\sigma=\sum_hp_h\sigma_h>0$.
\begin{enumerate}[label=(\roman*)]
  \item For all real $n_h>0$ with $\sum_hn_h=n$,
    \[
      \sum_h\frac{p_h^2\sigma_h^2}{n_h}\ge\frac{\bar\sigma^2}{n},
    \]
    with equality if and only if $n_h=np_h\sigma_h/\bar\sigma$ for every $h$. If every
    $\sigma_h>0$, this \emph{Neyman allocation} is the unique minimiser. If $\sigma_h=0$ for
    $k\ge1$ strata, the bound is not attained by an allocation with all $n_h>0$; one draw in
    each such stratum and the Neyman allocation of the other $n-k$ draws give
    $\bar\sigma^2/(n-k)$.
  \item $\displaystyle \frac{\sigma^2}{n}-\frac{\bar\sigma^2}{n}
    =\frac1n\sum_hp_h(\mu_h-\mu)^2+\frac1n\sum_hp_h(\sigma_h-\bar\sigma)^2$; the second term
    is the gain over proportional allocation, zero if and only if all $\sigma_h$ are equal.
  \item If the $n_h\ge1$ are computed from a pilot run independent of the draws used in
    \eqref{eq:st-estimator}, the estimator is unbiased, and conditionally on the pilot the
    allocation, though random, is fixed, so that \cref{prop:st-error}(i)--(ii) below apply.
\end{enumerate}
\end{proposition}

\begin{proof}
(i) By the Cauchy--Schwarz inequality,
\[
  \bar\sigma^2=\Big(\sum_h\frac{p_h\sigma_h}{\sqrt{n_h}}\sqrt{n_h}\Big)^2
  \le\Big(\sum_h\frac{p_h^2\sigma_h^2}{n_h}\Big)\Big(\sum_hn_h\Big),
\]
with equality if and only if the vectors $(p_h\sigma_h/\sqrt{n_h})_h$ and $(\sqrt{n_h})_h$ are
proportional, that is $n_h=cp_h\sigma_h$ for a constant $c$, which $\sum_hn_h=n$ fixes as
$n/\bar\sigma$. If some $\sigma_h=0$ this would require $n_h=0$. Strata with $\sigma_h=0$
contribute nothing to the variance whatever $n_h\ge1$, and the bound applied to the others
with $n-k$ draws gives $\bar\sigma^2/(n-k)$.

(ii) By \cref{thm:st-variance}(i),
$\sigma^2-\bar\sigma^2=\sum_hp_h(\mu_h-\mu)^2+\sum_hp_h\sigma_h^2-\bar\sigma^2$, and
$\sum_hp_h\sigma_h^2-\bar\sigma^2=\sum_hp_h(\sigma_h-\bar\sigma)^2$ is the variance of
$\sigma_h$ under the weights $p_h$.

(iii) Conditionally on the pilot run the $n_h$ are constants and the main draws keep their
law, so \cref{thm:st-variance}(ii) and \cref{prop:st-error} hold conditionally; for the mean,
take expectations.
\end{proof}

The $\sigma_h$ are estimated from a pilot run. If the pilot draws are reused in the final
estimate, the allocation depends on them, and both the estimate and its variance estimate are
in general biased (typically at order $1/n$, which we do not prove). Rounding to integers and
keeping $n_h\ge2$ costs little when $np_h\sigma_h/\bar\sigma\gg1$. For the test-case call with $10$ strata the Neyman
factor is $0.0234$ against $0.1007$ for proportional allocation (standard-error ratios
$0.153$ and $0.317$), and with $100$ strata $0.00044$ against $0.0081$. Half of the
$10$ strata lie below $z^*$; Neyman gives them nothing in the limit, and they could even be
skipped since $\mu_h=0$ is known. Neyman allocation is a form of importance sampling
(\cref{ch:09-importance-sampling}) restricted to a partition: it oversamples the strata where
the payoff varies, and the weights $p_hn/n_h$ implicit in \eqref{eq:st-estimator} are its
likelihood ratios, constant on each stratum. With unequal costs per stratum the rule changes
(\exref{ex:st-costs}).

\subsection{Error estimates}

\begin{proposition}[Standard errors for stratified estimators]\label{prop:st-error}
\begin{enumerate}[label=(\roman*)]
  \item If $n_h\ge2$ for every $h$ and $\hat s_h^2$ is the sample variance in stratum $h$,
    then $\hat v=\sum_hp_h^2\hat s_h^2/n_h$ is an unbiased estimator of
    $\Var\hat\mu^{\mathrm{st}}_n$.
  \item Let $H$ be fixed and $n_h/n\to q_h>0$ as $n\to\infty$, and let
    $v=\sum_hp_h^2\sigma_h^2/q_h>0$. Then $\sqrt n(\hat\mu^{\mathrm{st}}_n-\mu)\Rightarrow\mathcal N(0,v)$
    and $(\hat\mu^{\mathrm{st}}_n-\mu)/\sqrt{\hat v}\Rightarrow\mathcal N(0,1)$.
  \item $R$ independent replications of any stratified estimator (for instance with one draw
    per stratum) are i.i.d.\ with mean $\mu$, and \cref{thm:mc-clt} applies to their average
    as $R\to\infty$.
\end{enumerate}
\end{proposition}

\begin{proof}
(i) $\E\hat s_h^2=\sigma_h^2$ for $n_h\ge2$. (ii) By the central limit theorem in each
stratum, independently across $h$, $\sqrt{n_h}(\bar Y_h-\mu_h)\Rightarrow\mathcal N(0,\sigma_h^2)$,
a point mass if $\sigma_h=0$. As
$\sqrt n(\hat\mu^{\mathrm{st}}_n-\mu)=\sum_hp_h\sqrt{n/n_h}\,\sqrt{n_h}(\bar Y_h-\mu_h)$, the
characteristic functions multiply and the limit is $\mathcal N(0,v)$. By the strong law
$n\hat v\to v$ almost surely; Slutsky's lemma \citep[Lemma~2.8]{vandervaart-asymptotic}
concludes. (iii) is immediate.
\end{proof}

With one draw per stratum no $\hat s_h$ exists; (iii) is then the way to a confidence
interval, at the price of fewer strata per replication. Merging pairs of adjacent strata
gives a conservative alternative (\exref{ex:st-collapsed}).

\demo{10-stratification.ipynb}{stratify a call or a put at any strike with any number of
strata, compare proportional and Neyman allocation, and see where the draws go.}

\section{Post-stratification}\label{sec:st-post}

Sometimes the stratum of a draw can be observed but not prescribed: the paths are simulated
forward, and only afterwards do we see in which range their terminal value fell. If the
$p_h$ are known, the draws can still be reweighted. Let $X_1,\dots,X_n$ be i.i.d.\ copies of
$(\xi,Y)$, $N_h$ the number of draws with $\xi\in A_h$, $\bar Y_h$ the mean of their $Y$, and
$\mathcal E=\{N_h\ge1\text{ for every }h\}$. The \emph{post-stratified estimator} is
\[
  \hat\mu^{\mathrm{post}}_n=\sum_hp_h\bar Y_h\ \text{ on }\mathcal E,\qquad
  \hat\mu^{\mathrm{post}}_n=\hat\mu_n\ \text{ (the plain mean) on }\mathcal E^c .
\]

\begin{lemma}\label{lem:st-binomial}
If $N$ is binomial with parameters $n$ and $p\in(0,1]$, then
\[
  \E\big[N^{-1}\1_{\{N\ge1\}}\big]=\frac1{np}+\frac{1-p}{n^2p^2}+O(n^{-3}).
\]
\end{lemma}

\begin{proof}
Let $q=1-p$ and $j\ge1$. Since $\binom nk\frac{k!}{(k+j)!}=\frac{n!}{(n+j)!}\binom{n+j}{k+j}$,
the binomial theorem gives
\[
  \E\Big[\frac{1}{(N+1)\cdots(N+j)}\Big]=\frac{1-R_j}{(n+1)\cdots(n+j)\,p^j},
\]
where $R_j$ is the probability that a binomial$(n+j,p)$ variable is below $j$, so
$0\le R_j\le(n+j)^jq^{n}$. For $N\ge1$,
\[
  \frac1N=\frac1{N+1}+\frac1{(N+1)(N+2)}+\frac2{N(N+1)(N+2)},
\]
and $0\le\frac2{N(N+1)(N+2)}\le\frac8{(N+1)(N+2)(N+3)}$, whose expectation is $O(n^{-3})$.
Restricting the first two expectations to $\{N\ge1\}$ changes them by at most $q^n$. Hence
$\E[N^{-1}\1_{\{N\ge1\}}]=\frac1{(n+1)p}+\frac1{(n+1)(n+2)p^2}+O(n^{-3})$; finally
$\frac1{(n+1)p}=\frac1{np}-\frac1{n^2p}+O(n^{-3})$ and
$\frac1{(n+1)(n+2)p^2}=\frac1{n^2p^2}+O(n^{-3})$.
\end{proof}

\begin{proposition}[Post-stratification]\label{prop:st-post}
In this setting, with every $p_h>0$:
\begin{enumerate}[label=(\roman*)]
  \item Given the stratum labels of the $n$ draws, on $\mathcal E$ the estimator has
    conditional mean $\mu$ and conditional variance $\sum_hp_h^2\sigma_h^2/N_h$.
  \item $|\E\hat\mu^{\mathrm{post}}_n-\mu|\le\frac{\sigma}{\sqrt n}\big(\sum_h(1-p_h)^n\big)^{1/2}$:
    the bias is exponentially small.
  \item $\displaystyle\E\big[(\hat\mu^{\mathrm{post}}_n-\mu)^2\big]
    =\frac1n\sum_hp_h\sigma_h^2+\frac1{n^2}\sum_h(1-p_h)\sigma_h^2+O(n^{-3})$.
  \item $\sqrt n(\hat\mu^{\mathrm{post}}_n-\mu)\Rightarrow\mathcal N\big(0,\sum_hp_h\sigma_h^2\big)$.
\end{enumerate}
\end{proposition}

\begin{proof}
(i) The pairs (label, draw) are i.i.d., and given its label $h$ a draw has the conditional
law of $(\xi,Y)$ given $\xi\in A_h$; so, given all labels, the draws are independent with
these laws. On $\mathcal E$, $\bar Y_h$ is then a mean of $N_h$ independent draws from
stratum $h$; apply \cref{thm:st-variance}(ii) with $n_h=N_h$.

(ii) $\mathcal E$ is determined by the labels, so by (i)
$\E[(\hat\mu^{\mathrm{post}}_n-\mu)\1_{\mathcal E}]=0$, and
$|\E\hat\mu^{\mathrm{post}}_n-\mu|=|\E[(\hat\mu_n-\mu)\1_{\mathcal E^c}]|\le\frac{\sigma}{\sqrt n}\Prob(\mathcal E^c)^{1/2}$
by Cauchy--Schwarz, with $\Prob(\mathcal E^c)\le\sum_h(1-p_h)^n$.

(iii) By (i),
$\E[(\hat\mu^{\mathrm{post}}_n-\mu)^2]=\sum_hp_h^2\sigma_h^2\E[N_h^{-1}\1_{\mathcal E}]
+\E[(\hat\mu_n-\mu)^2\1_{\mathcal E^c}]$. Given the labels, the plain mean has conditional
mean $\sum_h\frac{N_h}n\mu_h$ and conditional variance $\sum_hN_h\sigma_h^2/n^2$, so
$\E[(\hat\mu_n-\mu)^2\mid\text{labels}]\le\max_h\sigma_h^2+\max_h(\mu_h-\mu)^2$, and the last
term is $O(\Prob(\mathcal E^c))$. Also
$0\le\E[N_h^{-1}\1_{\{N_h\ge1\}}]-\E[N_h^{-1}\1_{\mathcal E}]\le\Prob(\mathcal E^c)$. As
$\Prob(\mathcal E^c)\le H\max_h(1-p_h)^n=O(n^{-3})$, \cref{lem:st-binomial} with $N_h$
binomial$(n,p_h)$ gives
$\sum_hp_h^2\sigma_h^2\big(\frac1{np_h}+\frac{1-p_h}{n^2p_h^2}\big)+O(n^{-3})$, which is the claim.

(iv) On $\mathcal E$, whose probability tends to $1$,
$\hat\mu^{\mathrm{post}}_n=\varphi(\bar a,\bar b)$ with $\varphi(a,b)=\sum_hp_ha_h/b_h$,
$\bar a_h=\frac1n\sum_iY_i\1_{A_h}(\xi_i)$ and $\bar b_h=N_h/n$. The vectors
$(Y\1_{A_h}(\xi),\1_{A_h}(\xi))_h$ are i.i.d.\ and square integrable, so by the multivariate
central limit theorem and the delta method \citep[Theorem~3.1]{vandervaart-asymptotic},
$\sqrt n(\varphi(\bar a,\bar b)-\mu)$ converges to a centred normal law with the variance of
$\sum_h(\partial_{a_h}\varphi\,Y\1_{A_h}(\xi)+\partial_{b_h}\varphi\,\1_{A_h}(\xi))$ at
$(a_h,b_h)=(p_h\mu_h,p_h)$, where $\partial_{a_h}\varphi=1$ and $\partial_{b_h}\varphi=-\mu_h$.
That variable is $\sum_h\1_{A_h}(\xi)(Y-\mu_h)$, whose variance is $\sum_hp_h\sigma_h^2$.
\end{proof}

Post-stratification is therefore asymptotically as good as proportional stratification;
it loses the term $\frac1{n^2}\sum_h(1-p_h)\sigma_h^2$, the price of the random counts,
which is negligible when every $np_h$ is large. It breaks down when strata are so many that
some are often empty. It needs no conditional sampling, so it keeps whatever fast generator
the plain estimator uses. It is also a control variate in disguise, with the stratum
indicators as controls (\exref{ex:st-post-regression}). For the test-case call with $10$ strata and
$n=2\times10^5$, (iii) gives the standard-error ratio $0.3174$, larger than for proportional
allocation only in the sixth digit; the measured $0.311$ (\cref{tab:st-lab}) is below both by
sampling noise.

\section{Stratifying paths}\label{sec:st-paths}

A path-dependent payoff depends on many normals, $Y=f(Z)$ with $Z\sim\mathcal N(0,I_d)$.
Stratifying all of them on a product grid needs $H^d$ cells, which is hopeless beyond a few
dimensions. Two ways around this: stratify one well-chosen direction, or stratify every
coordinate separately (\cref{sec:st-lhs}).

\begin{proposition}[Stratifying a direction; terminal stratification]\label{prop:st-bridge}
\begin{enumerate}[label=(\roman*)]
  \item Let $Z\sim\mathcal N(0,I_d)$, $v\in\R^d$ a unit vector and $\xi=v^\top Z$. Then
    $\xi\sim\mathcal N(0,1)$, and $Z-v\xi$ is independent of $\xi$ with law
    $\mathcal N(0,I_d-vv^\top)$. Consequently, if $\xi'$ has the conditional law of $\xi$
    given $\xi\in A$ and $Z'\sim\mathcal N(0,I_d)$ is independent of $\xi'$, then
    $v\xi'+(I_d-vv^\top)Z'$ has the conditional law of $Z$ given $v^\top Z\in A$.
  \item Let $0<t_1<\dots<t_m=T$. In the Brownian-bridge construction of
    \cref{prop:ps-bridge}(iii), fill $t_m=T$ first, so that the first normal is $Z_1=W_T/\sqrt T$.
    If $Z_1$ is drawn from its conditional law given $Z_1\in A$ and $Z_2,\dots,Z_m$ are
    independent standard normals, the construction produces a vector with the conditional law
    of $(W_{t_1},\dots,W_{t_m})$ given $W_T/\sqrt T\in A$. With the increments construction,
    $W_T/\sqrt T=v^\top Z$ with $v_j=\sqrt{(t_j-t_{j-1})/T}$, and (i) gives the same law.
  \item Hence, for $f(W_{t_1},\dots,W_{t_m})\in L^2$ and $H$ equiprobable strata of $W_T$ with
    proportional allocation, the stratified estimator is unbiased with variance
    $\frac1n\sum_hp_h\sigma_h^2$, and $n\Var\hat\mu^{\mathrm{st}}_n\to\E[\Var(f\mid W_T)]$
    as $H\to\infty$.
\end{enumerate}
\end{proposition}

\begin{proof}
(i) $(\xi,Z-v\xi)$ is a linear image of $Z$, hence Gaussian, and
$\Cov(Z-v\xi,\xi)=v-v\,v^\top v=0$, so the two are independent; $\Var\xi=v^\top v=1$ and
$\Cov(Z-v\xi)=I_d-2vv^\top+vv^\top=I_d-vv^\top$. As $I_d-vv^\top$ is a symmetric projection,
$(I_d-vv^\top)Z'\sim\mathcal N(0,I_d-vv^\top)$. For bounded measurable $F$, independence gives
$\E[F(Z)\1_A(\xi)]=\E[F(v\xi+R)\1_A(\xi)]$ with $R=Z-v\xi$ independent of $\xi$; dividing by
$\Prob(\xi\in A)$ this is $\E[F(v\xi'+R')]$ with $\xi'$ and $R'\eqd R$ independent.

(ii) The construction is a fixed linear map, $W=AZ$, and $Z_1=W_T/\sqrt T$ because the
first step sets $W_T=\sqrt T Z_1$. The vector $(Z_1,\dots,Z_m)$ has independent components,
so conditioning on $Z_1\in A$ changes the law of $Z_1$ only; apply the map $A$. For the
increments construction $W_T=\sum_j\sqrt{t_j-t_{j-1}}Z_j$.

(iii) follows from (ii), \cref{thm:st-variance}(iii) and (v), with $\xi=N(W_T/\sqrt T)$.
\end{proof}

For geometric Brownian motion with deterministic parameters, the paths of
\cref{prop:ps-gbm-exact} built by the bridge on the variance clock have $\ln S_T$ set by
the first normal, so terminal stratification stratifies $S_T$; for a European payoff the
limit in (iii) is $0$. In the code, the function \texttt{terminal\_stratified\_normals}
draws $Z_1$ in its strata and the rest freely, and \texttt{paths.gbm\_paths} with
\texttt{method="bridge"} builds the paths (\cref{fig:st-paths}a).

\begin{figure}[ht]
  \centering
  \includegraphics[width=\textwidth]{ch10-paths}
  \caption{(a) Sixteen Brownian paths on $64$ steps, one in each of $16$ equiprobable
    strata of $W_1$ (ticks: stratum boundaries), completed by the Brownian bridge.
    (b) The arithmetic Asian call on $12$ monthly fixings in the test case: standard-error
    ratio to plain sampling at equal $n$ against the number $H$ of equiprobable strata, with
    proportional allocation and $n=2^{17}$ paths per point (seed $1021$). Stratified variables: $W_T$ through
    the bridge (blue), the first increment $Z_1$ of the increments construction (gold), and
    the direction of $\sum_jW_{t_j}$ (\cref{prop:st-bridge}(i); green). Horizontal lines:
    Latin hypercube sampling with the two constructions, from $2000$ independent designs of
    $2^{10}$ points.}
  \label{fig:st-paths}
\end{figure}

\begin{example}[An Asian call]\label{exa:st-asian}
The course's Asian example averages the stock at $t_j=j/12$, $j=1,\dots,12$, with strike
$100$ in the test case. The course reports $4.56058$, a Monte Carlo estimate from $10^8$
paths (standard error about $0.0008$); our geometric-control-variate estimate
(\texttt{asian\_reference}, two runs of $2^{25}$ paths, seeds $2026$ and $7$) is $4.56122$ and
$4.56127$, each with standard error $0.00003$, against the exact value $4.5612893$ computed in
\cref{exa:as-reference} (the first run is $2.0$ standard errors low); the course's estimate is
within its standard error. The plain standard deviation per path is $7.49$.
With $64$ strata of $W_T$ the standard-error ratio is $0.530$ (\cref{tab:st-lab}): terminal stratification removes about $72\%$ of the
variance of the payoff, and more strata add little (\cref{fig:st-paths}b), so this is nearly
all of the part that $W_T$ explains. Stratifying the first normal of the
increments construction instead fixes only $W_{1/12}$, and the ratio is $0.92$. The average
of the path matters more than its end: stratifying the direction $v\propto L^\top\mathbf 1$ of
$\sum_jW_{t_j}$, where $L$ is the increments matrix, gives $0.108$ with $64$ strata and about
$0.055$ with $256$ (\cref{fig:st-paths}b). That direction is, to within a cosine of $0.9999$,
the optimal importance-sampling drift of the payoff (\cref{prop:is-optimal-drift}), the
direction along which \textcite{glasserman-heidelberger-shahabuddin} stratify.
\end{example}

\subsection{Latin hypercube sampling}\label{sec:st-lhs}

Instead of one direction finely, Latin hypercube sampling stratifies every coordinate,
each into $n$ strata, with $n$ points in total.

\begin{definition}[Latin hypercube sample]\label{def:st-lhs}
Let $\pi_1,\dots,\pi_d$ be independent uniformly random permutations of $\{1,\dots,n\}$ and
$U_{ij}$, $i\le n$, $j\le d$, i.i.d.\ uniforms on $(0,1)$, independent of the permutations. The
points $V_i=(V_{i1},\dots,V_{id})$ with $V_{ij}=(\pi_j(i)-1+U_{ij})/n$ form a \emph{Latin
hypercube sample}, and $\hat\mu^{\mathrm{LHS}}_n=\frac1n\sum_if(V_i)$ estimates
$\mu=\int_{[0,1]^d}f$. Normals are obtained as $N^{-1}(V_{ij})$.
\end{definition}

Every column has exactly one point in each interval $[(k-1)/n,k/n)$, and the permutations
pair the columns at random (\cref{fig:st-lhs}a). For $f\in L^2([0,1]^d)$ write
$f_j(u_j)=\E[f(U)\mid U_j=u_j]-\mu$ for $U$ uniform on the cube, the \emph{additive part}
$f_{\mathrm{add}}=\mu+\sum_jf_j$ and the residual $e=f-f_{\mathrm{add}}$. The $f_j(U_j)$ are
independent and centred, and $e(U)$ is uncorrelated with each of them, so
$\sigma^2=\sum_j\Var f_j(U_j)+\int e^2$.

\begin{proposition}[Latin hypercube sampling]\label{prop:st-lhs}
Let $f\in L^2([0,1]^d)$ with variance $\sigma^2$ under the uniform law.
\begin{enumerate}[label=(\roman*)]
  \item Each $V_i$ is uniform on $[0,1]^d$, so $\hat\mu^{\mathrm{LHS}}_n$ is unbiased.
  \item If $f$ is additive, $f=\mu+\sum_jf_j$, then
    $\Var\hat\mu^{\mathrm{LHS}}_n=\sum_{j=1}^d\frac1{n^2}\sum_{k=1}^n\sigma_{j,k}^2$, with
    $\sigma_{j,k}^2$ the variance of $f_j(U)$ given $U\in[(k-1)/n,k/n)$: each additive term is
    integrated as by one-dimensional stratification with one point per stratum, and
    $n\Var\hat\mu^{\mathrm{LHS}}_n\to0$.
  \item \citep{stein-lhs} For every $f\in L^2$,
    $\Var\hat\mu^{\mathrm{LHS}}_n=\frac1n\int e^2+o(1/n)$: asymptotically, Latin hypercube
    sampling removes the variance of the additive part.
  \item \citep{mckay-beckman-conover} If $f$ is monotone in each coordinate,
    $\Var\hat\mu^{\mathrm{LHS}}_n\le\sigma^2/n$ for every $n$.
    \citep{owen-scrambled-nets} For every $f\in L^2$ and $n\ge2$,
    $\Var\hat\mu^{\mathrm{LHS}}_n\le\sigma^2/(n-1)$.
\end{enumerate}
\end{proposition}

\begin{proof}
(i) For fixed $i$, $\pi_j(i)$ is uniform on $\{1,\dots,n\}$ and independent of $U_{ij}$, so
$\Prob(V_{ij}\le x)=\frac1n\sum_k\Prob(k-1+U\le nx)=x$; the pairs $(\pi_j(i),U_{ij})$ are
independent across $j$.

(ii) Fix $j$ and list the points of column $j$ by stratum: the point in
$[(k-1)/n,k/n)$ is $(k-1+U_{\pi_j^{-1}(k)j})/n$. Given $\pi_j$, the $U_{\pi_j^{-1}(k)j}$,
$k=1,\dots,n$, are i.i.d.\ uniforms, and their conditional law does not depend on $\pi_j$. So
$\frac1n\sum_if_j(V_{ij})$ is the stratified estimator of $\int f_j$ with $n$ equiprobable
strata and one point each, whose variance is $\frac1{n^2}\sum_k\sigma_{j,k}^2$ by
\cref{thm:st-variance}(ii). Different columns are built from independent ingredients, so
the variances add. Finally
$n\Var\hat\mu^{\mathrm{LHS}}_n=\sum_j\sum_k\frac1n\sigma_{j,k}^2\to0$ by
\cref{thm:st-variance}(v) applied to $Y=f_j(U)$, for which $\Var(Y\mid U)=0$.

(iii) is \textcite{stein-lhs}; see also \citet[Proposition~10.1]{owen-mcbook}, where it
is stated for $\int f^2<\infty$. (iv) The first statement is the theorem stated in
Section~2.1 of \textcite{mckay-beckman-conover} and proved in their appendix (Section~8.3).
The second is due to \textcite{owen-scrambled-nets}, see \citet[Proposition~10.4]{owen-mcbook},
stated there for $n>2$; for $n=2$ it is immediate, since
$\Var\frac{f(V_1)+f(V_2)}2\le\frac12(\Var f(V_1)+\Var f(V_2))=\sigma^2$ by Cauchy--Schwarz.
\end{proof}

The $n$ points of a design are dependent, so the variance is estimated from independent
replications of the whole design, as in \cref{prop:st-error}(iii). The bound
$\sigma^2/(n-1)$ cannot be improved to $\sigma^2/n$: for a function with no additive part
Latin hypercube sampling can be slightly worse than plain sampling
(\exref{ex:st-lhs-interaction}). For the Asian call of \cref{exa:st-asian}, Latin hypercube
sampling of the $12$ normals gives standard-error ratios of $0.54\pm0.01$ with the increments
construction and $0.354\pm0.006$ with the bridge ($2000$ designs of $2^{10}$ points), which puts the most important variables first and
makes the payoff more nearly additive in them (\cref{tab:st-lab}). Latin hypercube sampling
balances only the one-dimensional projections of the design; the randomised quasi-Monte Carlo
point sets of \cref{ch:11-quasi-monte-carlo} also balance the higher-dimensional ones.

\begin{figure}[ht]
  \centering
  \includegraphics[width=\textwidth]{ch10-lhs}
  \caption{(a) A Latin hypercube sample of $10$ points in $[0,1]^2$: one point in each row
    and each column of the grid. (b) $n\Var\hat\mu^{\mathrm{LHS}}_n/\sigma^2$ against $n$ in
    dimension $5$, from $400$ independent designs for each $n$ (seed $1040$), for $\sum_je^{u_j}$ (additive),
    the same plus $6(u_1-\frac12)(u_2-\frac12)$, and $\prod_j2u_j$. Dashed: the limits
    $\int e^2/\sigma^2$ of \cref{prop:st-lhs}(iii), $0.17$ and $0.48$; the solid line is plain
    sampling.}
  \label{fig:st-lhs}
\end{figure}

\demo{10-stratification.ipynb}{stratify Brownian paths on their endpoint or on any linear
combination of the path, and draw Latin hypercube designs.}

\section{Conditional Monte Carlo}\label{sec:st-conditional}

If part of the randomness can be integrated out in closed form, it should be.

\begin{theorem}[Conditional Monte Carlo]\label{thm:st-conditional}
Let $Y\in L^2$, let $\mathcal G$ be a sub-$\sigma$-algebra, and $Y^{\mathrm c}=\E[Y\mid\mathcal G]$.
\begin{enumerate}[label=(\roman*)]
  \item $\E Y^{\mathrm c}=\E Y$ and $\Var Y=\Var Y^{\mathrm c}+\E[\Var(Y\mid\mathcal G)]$,
    where $\Var(Y\mid\mathcal G)=\E[(Y-Y^{\mathrm c})^2\mid\mathcal G]$. Hence
    $\Var Y^{\mathrm c}\le\Var Y$, with equality if and only if $Y=Y^{\mathrm c}$ almost
    surely, that is, if and only if $Y$ is almost surely equal to a $\mathcal G$-measurable variable.
  \item If $\mathcal G_1\subset\mathcal G_2$, then
    $\Var\E[Y\mid\mathcal G_1]\le\Var\E[Y\mid\mathcal G_2]$.
  \item The average $\hat\mu^{\mathrm c}_n$ of $n$ i.i.d.\ copies of $Y^{\mathrm c}$ is
    unbiased with variance $\Var Y^{\mathrm c}/n$; if a copy of $Y^{\mathrm c}$ costs $\lambda$
    times a copy of $Y$, its efficiency relative to plain sampling is
    $\Var Y/(\lambda\Var Y^{\mathrm c})$, which may be below $1$.
\end{enumerate}
\end{theorem}

\begin{proof}
(i) $\E Y^{\mathrm c}=\E Y$ by the tower property. $Y-Y^{\mathrm c}$ is orthogonal in $L^2$ to
every square-integrable $\mathcal G$-measurable variable, in particular to
$Y^{\mathrm c}-\E Y$, so
$\E(Y-\E Y)^2=\E(Y^{\mathrm c}-\E Y)^2+\E(Y-Y^{\mathrm c})^2$, and
$\E(Y-Y^{\mathrm c})^2=\E[\Var(Y\mid\mathcal G)]$. The last term vanishes if and only if
$Y=Y^{\mathrm c}$ almost surely. (ii) Apply (i) to $\E[Y\mid\mathcal G_2]$ and $\mathcal G_1$, using
$\E[\E[Y\mid\mathcal G_2]\mid\mathcal G_1]=\E[Y\mid\mathcal G_1]$. (iii) is
\cref{def:mc-efficiency}.
\end{proof}

In statistics this is the Rao--Blackwell theorem \citep[Theorem~7.3.17]{casella-berger},
where $\mathcal G$ is generated by a sufficient statistic so that $Y^{\mathrm c}$ does not
depend on the unknown parameter. Here $\mathcal G$ is ours to choose, and the only
requirement is that $Y^{\mathrm c}$ be computable. Part (ii) is the trade-off: conditioning
on less information removes more variance but leaves a harder conditional expectation.
\Cref{eq:st-total-variance} shows how conditioning and stratification divide the work. For
a scalar $\xi$ with a continuous law and many equiprobable strata, stratifying on $\xi$
gives $n\Var\to\E[\Var(Y\mid\xi)]$ (\cref{thm:st-variance}(v)), and conditioning on $\xi$ gives
$n\Var=\Var(\E[Y\mid\xi])$; the two add up to $\sigma^2$. Stratify on what cannot be integrated,
integrate out what can, and, when $\E[Y\mid\xi]=g(\xi)$ is available, do both: stratifying
$\xi$ in $g(\xi)$ has variance $o(1/n)$. Further examples in these notes are the
Brownian-bridge correction for barriers, which integrates out the path between monitoring
dates (\cref{prop:bar-bridge-crossing}), conditioning a two-asset payoff on one asset
(\exref{ex:st-spread}) and Curran's conditioning on the geometric average
(\cref{prop:as-curran}). Conditional values are smooth in the parameters even when the payoff
is not, which is what pathwise Greeks need (\cref{thm:gr-pathwise}).

\subsection{Conditioning on the volatility path}

In a stochastic volatility model, given the whole path of the variance the log-price is
Gaussian, and the Black--Scholes formula gives the conditional price
\citep{hull-white-sv,romano-touzi,willard}.

\begin{proposition}[Mixing: conditioning on the volatility path]\label{prop:st-mixing}
Let $W$ and $W^\perp$ be independent Brownian motions under $\Q$, $(\F^W_t)$ the
filtration generated by $W$, and $V\ge0$ an $(\F^W_t)$-progressively measurable process
with $Q_T=\int_0^TV_t\,\dd t<\infty$ almost surely. Let $\rho\in[-1,1]$, $\bar\rho=\sqrt{1-\rho^2}$,
$I_T=\int_0^T\sqrt{V_t}\,\dd W_t$, and
\[
  S_T=S\exp\Big((r-y)T-\tfrac12Q_T+\rho I_T+\bar\rho\int_0^T\sqrt{V_t}\,\dd W^\perp_t\Big),
\]
the solution of $\dd S_t=S_t\big((r-y)\,\dd t+\sqrt{V_t}(\rho\,\dd W_t+\bar\rho\,\dd W^\perp_t)\big)$. Set
\[
  \tilde S=S\exp\big(\rho I_T-\tfrac12\rho^2Q_T\big),\qquad \tilde\sigma=\bar\rho\sqrt{Q_T/T}.
\]
\begin{enumerate}[label=(\roman*)]
  \item There is a standard normal $N$, independent of $\F^W_T$, with
    $\ln S_T=\ln\tilde S+(r-y-\frac12\tilde\sigma^2)T+\tilde\sigma\sqrt TN$.
  \item For every Borel $f$ with $\E|f(S_T)|<\infty$,
    $\E[f(S_T)\mid\F^W_T]=\int_\R f\big(\tilde Se^{(r-y-\tilde\sigma^2/2)T+\tilde\sigma\sqrt Tz}\big)N'(z)\,\dd z$.
    In particular the conditional price of a call is the Black--Scholes price at spot
    $\tilde S$ and volatility $\tilde\sigma$, $e^{-rT}\E[(S_T-K)^+\mid\F^W_T]=C_{\mathrm{BS}}(\tilde S,K,T,r,y,\tilde\sigma)$,
    and likewise for a put; when $\tilde\sigma=0$ (for $\rho=\pm1$ or $Q_T=0$) it is the
    discounted payoff at the forward $\tilde Se^{(r-y)T}$.
  \item (\emph{Hull and White}.) If $\rho=0$, the call price is
    $\E\big[C_{\mathrm{BS}}(S,K,T,r,y,\sqrt{Q_T/T})\big]$: a mixture of Black--Scholes prices
    over the law of the realised volatility (\exref{ex:st-hull-white} expands it for small
    volatility of variance).
  \item (\emph{On a grid}.) If $\ln S$ is advanced by the log-Euler step
    $\Delta\ln S=(r-y-\frac12V_{t_k})\Delta_k+\sqrt{V_{t_k}}(\rho\Delta W_k+\bar\rho\Delta W^\perp_k)$
    with $V_{t_k}$ functions of $W$ only, (i) and (ii) hold exactly with $Q_T$ and $I_T$ replaced
    by $\sum_kV_{t_k}\Delta_k$ and $\sum_k\sqrt{V_{t_k}}\Delta W_k$.
\end{enumerate}
\end{proposition}

\begin{proof}
(i) Let $\mathcal G_t=\F^W_T\vee\sigma(W^\perp_s:s\le t)$, completed. Since $W^\perp$ is
independent of $W$ and has independent increments, $W^\perp_t-W^\perp_s$ is independent of
$\mathcal G_s$, so $W^\perp$ is a $(\mathcal G_t)$-Brownian motion. Let $h_t=\sqrt{V_t/Q_T}$
on $\{Q_T>0\}$ and $h_t=1/\sqrt T$ on $\{Q_T=0\}$. It is measurable in $(t,\omega)$ with
respect to $\mathcal B([0,T])\otimes\F^W_T$, hence $(\mathcal G_t)$-progressive, and
$\int_0^Th_t^2\,\dd t=1$. Let $N=\int_0^Th\,\dd W^\perp$ in the filtration $(\mathcal G_t)$.
For $u\in\R$, It\^o's formula shows that
$M_t=\exp\big(iu\int_0^th\,\dd W^\perp+\frac{u^2}2\int_0^th^2\,\dd s\big)$ is a local martingale;
it is bounded by $e^{u^2/2}$, hence a martingale, and $\E[M_T\mid\mathcal G_0]=1$, that is,
$\E[e^{iuN}\mid\mathcal G_0]=e^{-u^2/2}$. A deterministic conditional characteristic function
means that $N$ is standard normal and independent of $\mathcal G_0\supset\F^W_T$. Finally,
$\int_0^T\sqrt V\,\dd W^\perp=\sqrt{Q_T}\,N$ almost surely: $\sqrt V=\sqrt{Q_T}\,h$, the
$\mathcal G_0$-measurable factor $\sqrt{Q_T}$ can be taken out of the integral (check it on
simple integrands), and the integral of the $(\F^W_t\vee\F^{W^\perp}_t)$-progressive
process $\sqrt V$ is the same in either filtration, since $W^\perp$ is a Brownian motion for
both and the integral is the limit in probability of the integrals of the same simple
processes. Substitute in $\ln S_T$: $\bar\rho\sqrt{Q_T}=\tilde\sigma\sqrt T$ and
$-\frac12Q_T+\rho I_T=\ln(\tilde S/S)-\frac12\tilde\sigma^2T$.

(ii) $\tilde S$ and $\tilde\sigma$ are $\F^W_T$-measurable and $N$ is independent of
$\F^W_T$; apply the independence lemma \citep[Lemma~2.3.4]{shreve2}. For the call, the
integral is the Black--Scholes expectation with spot $\tilde S$ and volatility
$\tilde\sigma$. (iii) is (ii) with $\tilde S=S$ and $\tilde\sigma=\sqrt{Q_T/T}$, after taking
expectations. (iv) Given $\F^W_T$ the sum $\sum_k\sqrt{V_{t_k}}\Delta W^\perp_k$ is a
combination, with $\F^W_T$-measurable coefficients, of the independent
$\Delta W^\perp_k\sim\mathcal N(0,\Delta_k)$, hence conditionally $\mathcal N(0,\sum_kV_{t_k}\Delta_k)$; the
rest is as before.
\end{proof}

The conditional estimator averages $C_{\mathrm{BS}}(\tilde S,K,T,r,y,\tilde\sigma)$ over
simulated variance paths and needs no normal for $W^\perp$. By \cref{thm:st-conditional} it
removes the part of the variance due to $W^\perp$, which is large when $|\rho|$ is small. At
$\rho=\pm1$ there is nothing left to integrate. In the Heston model (\cref{def:sv-heston}),
$\dd V_t=\kappa(\theta-V_t)\,\dd t+\xi\sqrt{V_t}\,\dd W_t$; this equation has a pathwise
unique strong solution, adapted to the filtration of $W$ \citep[Sections~5.2 and 5.3]{karatzas-shreve},
and integrating it gives
\[
  I_T=\frac1\xi\big(V_T-V_0-\kappa\theta T+\kappa Q_T\big),
\]
so the conditional estimator needs only $V_T$ and $Q_T$ \citep{willard,broadie-kaya}.

\paragraph{A laboratory.} Take the test case with $V_0=\theta=0.04$, $\kappa=2$ and
$\xi=0.3$, and the at-the-money put, whose payoff is bounded, so that every variance is
finite whatever the moments of $S_T$. We sample $V$ exactly at $64$ equal steps, by
noncentral chi-square transitions \citep[Section~3.4.1]{glasserman-mc}, take $Q_T$ by the
trapezoidal rule and $I_T$ from the identity above. The plain estimator draws one more normal
per path and builds $S_T$ from (i); the conditional estimator uses (ii) on the same variance
paths. Both estimate the same discretised price, and the conditional one is the
Rao--Blackwellisation of the other. The Fourier prices of \cref{prop:sv-fourier} are $8.1534$
for $\rho=0$ and $7.9480$ for $\rho=-0.7$. Over $40$ independent runs of $5000$ paths
(our implementation, numpy \texttt{default\_rng}, seeds $777$ to $816$), the errors of the
conditional estimate divided by its standard error have mean $0.16$ and standard deviation
$1.04$ for $\rho=0$, and $0.15$ and $1.01$ for $\rho=-0.7$: no discretisation bias is visible,
and the standard errors are honest. The standard-error ratio of the conditional to the plain estimator is
$0.163$ for $\rho=0$ and $0.728$ for $\rho=-0.7$ ($4\times10^5$ paths, \cref{tab:st-lab}); the
cost per path is dominated by the $64$ variance steps and is the same for both, so the
efficiencies are about $37$ and $1.9$. \Cref{fig:st-mixing} shows why: for $\rho=0$ the
conditional value is a smooth function of the realised variance, while the plain payoff
scatters around it.

\begin{figure}[ht]
  \centering
  \includegraphics[width=\textwidth]{ch10-mixing}
  \caption{Conditional Monte Carlo for the at-the-money put in the Heston model ($V_0=\theta=0.04$,
    $\kappa=2$, $64$ exact variance steps). (a) $\rho=0$, $\xi=0.3$: plain payoffs and conditional values
    $P_{\mathrm{BS}}(S,K,T,r,y,\sqrt{Q_T/T})$ for $2000$ variance paths, against the realised volatility.
    (b) Standard-error ratio of the conditional to the plain estimator against $-\rho$, from
    $5\times10^4$ variance paths (the same for every $\rho$; seed $1031$), for $\xi=0.3$ and $\xi=0.6$.}
  \label{fig:st-mixing}
\end{figure}

\demo{10-stratification.ipynb}{conditional Monte Carlo in the Heston model: move $\rho$ and
$\xi$ and watch the plain payoffs collapse onto the conditional values.}

\section{Laboratory}\label{sec:st-lab}

\Cref{tab:st-lab} collects the methods of the chapter. One-dimensional rows use $n=2\times10^5$
draws (the test-case call) and $n=10^5$ (the course's deep out-of-the-money call); the
exact ratios come from \cref{thm:st-variance,prop:st-optimal,prop:st-post} with stratum
moments by quadrature. The ratios are produced by \texttt{stratified.laboratory()} (numpy
\texttt{default\_rng}, seed $10$). The costs are medians of wall-clock times of the whole
estimator against its plain competitor, measured by \texttt{laboratory\_costs()} on our machine:
they vary by tens of percent from run to run and more across machines, and the notebook
measures them on yours.

\begin{table}[ht]
  \centering
  \small
  \begin{tabular}{@{}llcccc@{}}
    \toprule
    target & method & exact & measured & cost & efficiency\\
    \midrule
    call & $10$ strata, proportional & $0.317$ & $0.318$ & $4.5$ & $2.2$\\
     & $10$ strata, Neyman & $0.153$ & $0.152$ & $5$ & $8.6$\\
     & $10$ strata, post-stratified & $0.317$ & $0.311$ & $3$ & $3.3$\\
     & $100$ strata, proportional & $0.090$ & $0.085$ & $4.5$ & $27$\\
    \midrule
    deep OTM call & $10$ strata, proportional & $0.925$ & $0.939$ & $4.5$ & $0.26$\\
     & $100$ strata, proportional & $0.398$ & $0.386$ & $4.5$ & $1.4$\\
    \midrule
    Asian call & $64$ strata of $W_T$ (bridge) & & $0.530$ & $1.6$ & $2.2$\\
     & $64$ strata of the first increment & & $0.917$ & $1.2$ & $1.0$\\
     & $64$ post-strata of $W_T$ & & $0.530$ & $1.5$ & $2.4$\\
     & $64$ strata of $\sum_jW_{t_j}$ & & $0.108$ & $1.5$ & $57$\\
     & Latin hypercube, increments & & $0.54\pm0.01$ & $4.2$ & $0.8$\\
     & Latin hypercube, bridge & & $0.354\pm0.006$ & $4.8$ & $1.7$\\
    \midrule
    Heston put & conditional, $\rho=0$ & & $0.163$ & $1.0$ & $37$\\
     & conditional, $\rho=-0.7$ & & $0.728$ & $1.0$ & $1.9$\\
    \bottomrule
  \end{tabular}
  \caption{Standard-error ratios $\varrho$ to plain Monte Carlo at equal $n$, exact and measured;
    cost per draw relative to the plain estimator (NumPy normals; increments construction for the
    Asian call), on our machine; and the efficiency $1/(\varrho^2\times\text{cost})$, from the exact
    ratio where there is one. Measured ratios divide the method's standard error by
    $\sigma/\sqrt n$, with $\sigma$ exact for the calls, from $2^{20}$ plain paths for the Asian
    call, and from the plain estimator on the same paths for the Heston put. The Asian rows use
    $n=2^{16}$ paths, Latin hypercube $2000$ designs of $2^{10}$ points; the Heston rows use
    $4\times10^5$ variance paths, whose simulation dominates the cost of both estimators.}
  \label{tab:st-lab}
\end{table}

Costs change the picture. Stratification by inversion pays in one dimension only with enough
strata, and loses for the deep out-of-the-money call with $10$ strata, whose equiprobable strata
are wasted where the payoff vanishes; post-stratification keeps the fast generator. Latin
hypercube sampling needs $12$ inversions per path and barely covers that cost here, while
stratifying the right direction costs little and gains a factor $57$. Conditioning costs one
Black--Scholes evaluation per path and helps as much as the integrated-out noise contributes.

\paragraph{Course examples reproduced.} The stratified-sampling experiment of the course's
notebook, with its exact variances; the call price $7.36429$; and the course's estimate
$4.56058$ of its Asian call, confirmed within its standard error.

\furtherreading
\textcite{glasserman-mc} is the reference: Section~4.3.1 for stratified sampling and optimal
allocation, Section~4.3.2 for stratified paths, Section~4.3.3 for post-stratification and
Section~4.4 for Latin hypercube sampling; \textcite{asmussen-glynn}, Sections~V.4 and~V.7,
treat conditional Monte Carlo and stratification. \textcite{neyman-1934} introduced the optimal
allocation; \textcite{cochran}, Chapter~5, is the classical account of stratified sampling. Latin hypercube sampling is due to \textcite{mckay-beckman-conover}; its asymptotics
are in \textcite{stein-lhs} and \textcite{owen-lhs-clt}, and \textcite{owen-mcbook},
Chapter~10, gives a modern account. The mixing formula goes back to \textcite{hull-white-sv}
for independent volatility; \textcite{romano-touzi} treat correlated volatility, and
\textcite{willard} turned it into a conditional Monte Carlo method.
